\documentclass[a4paper,12pt]{extarticle}
\usepackage[margin=1.5cm]{geometry}
\usepackage{amsmath, amssymb}
\usepackage{array}
\usepackage{multicol}
\usepackage{helvet}
\usepackage{multirow}
\usepackage{graphicx}
\renewcommand{\baselinestretch}{2}
\renewcommand{\familydefault}{\sfdefault}

\begin{document}


    \begin{tabular}{| c | c | c | c |}
        \hline
        & MATEMÁTICAS & 27/02/2026 & calificación \\
        \hline 

        \multirow{2}{*}{\includegraphics{../../img-Logos/iesoGalileo.jpg}} & Nombre y apellidos \hspace{120pt} & Ex. 3ºESO-D UD5:  & \\ 
        &&Ecuaciones& \\ 
        \hline
    \end{tabular}
    
    \begin{tabular}{| l | }
        \hline
        NOTAS: \\ 
            1. El examen tiene que ser limpio, ordenado y sin faltas de ortografía. \\
            
            2. El examen ha de realizarse en bolígrafo negro o azul, evitando tachones en la medida \\
            de lo posible. \\ 
            
            3. Deben aparecer todas las operaciones NO vale con indicar solo el resultado.\\
            
            4. Los problemas deben contener: Datos, planteamiento y resolución, respondiendo a lo \\
            que se pregunte, NO vale con indicar un número como solución del problema.\\

            5. Persona que se vea copiando o hablando con un compañero tiene un cero en el examen\\
            
        \hline
    \end{tabular}


    \begin{enumerate}

        \item Dados los polinomios: 
        \[P(x) = 2x^3-4x^2+5x-6\]
        \[Q(x) = x^2 -3x+2\]
        \[R(x) = x-4\]

        \begin{enumerate}
            \item (0,25 puntos) Indica el grado, término principal y término independiente de P(x).
            \item (0,5 puntos) Halla el valor númerico de P(x) para x=2.
            \item (0,5 puntos) ¿Es x=1 raíz de alguno de los polinomios? ¿Cúal? ¿Por qué?
            \item (1 punto) Realiza las siguentes operaciones:
            \begin{itemize}
                \item (0,3p) $P-2Q$
                \item (0,3p) $ P \cdot Q $ 
                \item (0,4p) $ P : Q $ 
            \end{itemize} 
        \end{enumerate}

        \item (0,75 puntos) Factoriza el siguiente polinommio 
        \[P(x)= 3x^5 -9x^4-3x^3+9x^2\]
        \item (1p) Resuelve la siguiente ecuación:
        \[\dfrac{x}{2}-\dfrac{x+3}{4}=\dfrac{x-1}{3}\]
        \item (1,5 puntos) Resuelve la siguiente ecuación
        \[(2x+2)(2x-2)=(x+1)^2 +2(x+1)(x-1)\]
        \item (1,5 puntos) Halla dos números enteros consecutivos tales que la suma de sus cuadrados es 145. 
        \item (1,5 puntos) Resuelve el siguiente sistema: 
        $\left \{\begin{array}{rr}
                    3x+5y=1 \\
                    x+6y=-4
                \end{array}\right .$  

        \item (1,5 puntos) Tenemos 53 céntimos de euro repartidos en 16 monedas, unas de dos céntimos y otras de cinco céntimos. ¿Cuántas monedas de cada clase tenemos?
        
    \end{enumerate}


\end{document}